\section*{Введение}
\addcontentsline{toc}{section}{Введение}

Распознавание именованных сущностей (Named Entity Recognition, NER) --- базовая задача обработки естественного языка, которая состоит в выделении в тексте упоминаний объектов заданных категорий и их классификации. В научных текстах набор категорий существенно отличается от новостных или общелингвистических корпусов: вместо персон и организаций исследователям важны методы, метрики, датасеты и научные концепции.

Целью данного проекта является сравнительное исследование нейросетевых архитектур на задаче NER научных текстов. В качестве датасета использован SciERC~\cite{luan2018scierc} --- корпус из 500 аннотированных научных абстрактов по компьютерным наукам. Набор содержит 6 типов именованных сущностей, что делает его нетривиальным: классы семантически близки и имеют существенный дисбаланс по частоте.

В работе решаются следующие задачи:
\begin{enumerate}
    \item Реализация и обучение 10 моделей на базе предобученных трансформеров (BERT, RoBERTa, SciBERT) с различными классификационными головами.
    \item Исследование влияния взвешивания классов, архитектуры головы и качества предобучения энкодера на результат.
    \item Оценка влияния улучшенного препроцессинга данных: аугментации через замену сущностей и расширения контекста через соседние предложения.
\end{enumerate}
